%  CONFIGURE NEW SINGLE-PAGE FORMAT (from template 02_Article_Supplementary.tex)

\onecolumn % go back to one column
\fancyhead{} % make sure we get no headers
\renewcommand{\floatpagefraction}{0.1}
\lfoot[\bSupInf]{\dAuthor}
\rfoot[\dAuthor]{\cSupInf}
\newpage

\captionsetup*{format=largeformat} % make figure legend slightly larger than in the paper
\setcounter{figure}{0} % reset figure counter for Supp. Figures
\setcounter{table}{0} % reset table counter for Supp. Tables
\setcounter{equation}{0} % reset equation counter for Supp. Equations
\makeatletter
\renewcommand{\thefigure}{S\@arabic\c@figure} % make Figure legend start with Figure S
\renewcommand{\thetable}{S\@arabic\c@table} % make Table legend start with Table S
\makeatother
\def\theequation{S\arabic{equation}}

%  APPENDIX CONTENT (from 2_thesis/main.tex)

\appendix
\section{Figures and Tables}

\subsection{Detailed Experiment Configuration}
\begin{table*}[h]
\centering
\caption{Method configurations for the recovery-radius benchmark. The three gradient methods share an identical optimizer, surrogate, and parameter-transform setup; they differ only in the loss function. \texttt{nm\_guarino} optimizes the same soft Guarino loss with Nelder--Mead, evaluating but not consuming the surrogate gradient. \texttt{nm\_guarino\_hard} optimizes the same loss composition with \emph{hard} feature extraction (threshold-crossing spike detection on both simulated and target traces), so the soft-feature hyperparameters $\tau$, $\beta$, $\beta_v$ vanish from its configuration.}
\label{tab:training-params}
\footnotesize
\setlength{\tabcolsep}{3pt}
\begin{tabular}{lccccc}
\toprule
\textbf{Parameter} & \texttt{grad\_mse} & \texttt{grad\_guarino} & \texttt{grad\_vanrossum} & \texttt{nm\_guarino} & \texttt{nm\_guarino\_hard} \\
\midrule
\multicolumn{6}{l}{\textit{Optimizer}} \\
Algorithm                 & \multicolumn{3}{c}{Polyak SGD}                                & \multicolumn{2}{c}{Nelder--Mead} \\
Learning rate             & \multicolumn{3}{c}{$0.001$}                                   & \multicolumn{2}{c}{---} \\
Polyak $(\alpha,\beta)$   & \multicolumn{3}{c}{$(0.8,\,0.6)$}                             & \multicolumn{2}{c}{---} \\
Steps / max function evals & \multicolumn{3}{c}{25 gradient steps}                         & \multicolumn{2}{c}{50 evals} \\
NM tolerances $(x_{\text{atol}}, f_{\text{atol}})$ & \multicolumn{3}{c}{---}                & \multicolumn{2}{c}{$(10^{-4},\,10^{-4})$, adaptive} \\
Parameter transform       & \multicolumn{3}{c}{sigmoid (\cref{sec:method:opt:stabilization:sigmoid})} & \multicolumn{2}{c}{box-bounded} \\
\midrule
\multicolumn{6}{l}{\textit{Surrogate gradient (gradient methods only)}} \\
Surrogate type            & \multicolumn{3}{c}{SuperSpike (\cref{eq:superspike-surrogate})} & \multicolumn{2}{c}{not consumed} \\
Surrogate slope $\beta$   & \multicolumn{3}{c}{$6.0$}                                     & \multicolumn{2}{c}{---} \\
\midrule
\multicolumn{6}{l}{\textit{Loss-specific parameters}} \\
Feature extraction        & ---   & soft    & ---     & soft    & hard \\
MSE: normalize            & False & ---     & ---     & ---     & --- \\
MSE: voltage clamp        & none  & ---     & ---     & ---     & --- \\
Guarino: feature temp.\ $\tau$ & --- & $0.3$ & ---   & $0.3$   & --- \\
Guarino: sigmoid sharpness $\beta$ & --- & $10.0$ & --- & $10.0$ & --- \\
Guarino: validity $\beta_v$ & --- & $10.0$ & ---    & $10.0$  & --- \\
Guarino: spike-count weight & --- & $0.1$ & ---     & $0.0$   & $0.0$ \\
Guarino: missing-feature penalty & --- & $3.0$ & --- & $3.0$  & $3.0$ \\
Van Rossum: kernel $\tau$ & ---   & ---     & $8$\,ms & ---   & --- \\
Van Rossum: subthreshold $\lambda_{\text{sub}}$ & --- & --- & $0.3$ & --- & --- \\
\midrule
\multicolumn{6}{l}{\textit{Simulation (shared)}} \\
Time step $\Delta t$       & \multicolumn{5}{c}{$0.025$\,ms} \\
Total / stimulus duration & \multicolumn{5}{c}{$100\,$ms / $90\,$ms (5\,ms delay)} \\
Spike threshold            & \multicolumn{5}{c}{$0$\,mV} \\
\bottomrule
\end{tabular}
\end{table*}

\subsection{Loss Landscape Analysis}\label{sec:app:loss}
To characterize the optimization problem independently of any particular training run, we computed each of the four loss functions on a dense two-parameter grid ($20 \times 20 = 400$ points per pair) for 28 parameter pairs, for a total of $11{,}200$ simulations and $\sim$3.8\,h of compute.
The remaining parameters were held at the default initial values of \cref{tab:adex_params}.

\newgeometry{total={130mm,247mm},left=40mm,top=25mm} % thesis text width (130 mm) so the landscape-grid figures keep their thesis size
\onecolumn % re-apply single column after the geometry change

\begin{figure*}[h]
  \centering
  \begin{subfigure}{\textwidth}
    \includegraphics[width=\textwidth]{figures/mse_loss_landscape_grid.pdf}
    \caption{Loss landscape.}\label{fig:results:mse-loss-landscape}
  \end{subfigure}
  \begin{subfigure}{\textwidth}
    \includegraphics[width=\textwidth]{figures/mse_grad_landscape_grid.pdf}
    \caption{Gradient landscape.}\label{fig:results:mse-grad-landscape}
  \end{subfigure}
  \caption{The MSE loss landscape --- even in these simplified conditions where only two parameters at a time are changed ---
is characterized by lots of local minima. Some parameter combinations are also characterized by long shallow valleys.}
\end{figure*}

\begin{figure*}[h]
  \centering
  \begin{subfigure}{\textwidth}
    \includegraphics[width=\textwidth]{figures/soft-dtw_loss_landscape_grid.pdf}
    \caption{Loss landscape.}\label{fig:results:sdtw-loss-landscape}
  \end{subfigure}
  \begin{subfigure}{\textwidth}
    \includegraphics[width=\textwidth]{figures/soft-dtw_grad_landscape_grid.pdf}
    \caption{Gradient landscape.}\label{fig:results:sdtw-grad-landscape}
  \end{subfigure}
  \caption{When observing the \gls{SDTW} loss landscape, we identify very thin and deep valleys containing the global minimum.
This correlates with the observation, that \gls{SDTW} was only able to recover very small parameter perturbations.
For larger perturbations, gradients were close to 0, which can be explained by the large loss plateaus visible.}
\end{figure*}

\begin{figure*}[h]
  \centering
  \begin{subfigure}{\textwidth}
    \includegraphics[width=\textwidth]{figures/vanrossum_loss_landscape_grid.pdf}
    \caption{Loss landscape.}\label{fig:results:vanrossum-loss-landscape}
  \end{subfigure}
  \begin{subfigure}{\textwidth}
    \includegraphics[width=\textwidth]{figures/vanrossum_grad_landscape_grid.pdf}
    \caption{Gradient landscape.}\label{fig:results:vanrossum-grad-landscape}
  \end{subfigure}
  \caption{The Van Rossum loss promises to be a more interesting candidate based on these loss landscape plots.
For example, the local minima are less pronounced than in the \gls{MSE} landscape, even though still present.
The global-minimum valleys tend to be less narrow compared to \gls{SDTW}; however, they still tend to be surrounded by large loss plateaus.}
\end{figure*}

\begin{figure*}[h]
  \centering
  \begin{subfigure}{\textwidth}
    \includegraphics[width=\textwidth]{figures/guarino_loss_landscape_grid.pdf}
    \caption{Loss landscape.}\label{fig:results:guarino-loss-landscape}
  \end{subfigure}
  \begin{subfigure}{\textwidth}
    \includegraphics[width=\textwidth]{figures/guarino_grad_landscape_grid.pdf}
    \caption{Gradient landscape.}\label{fig:results:guarino-grad-landscape}
  \end{subfigure}
  \caption{The feature-based loss shares lots of similarities with the Van-Rossum loss.
It seems to be the most resistant against local minima in this testing configuration.
However, it shares the large loss plateaus which could indicate small-to-zero gradients for ill conditioned initial parameters.}
\end{figure*}

\clearpage
\subsection{Additional Results Figures}
Supporting figures for \cref{ch:results}.

\begin{figure}[H]
  \begin{center}
    \includegraphics[width=\textwidth]{figures/surrogate_vs_non-surrogate_comparison}
  \end{center}
  \caption{Comparison of \texttt{AdEx} (black, solid) and \texttt{AdExSurrogate} (red, dashed; sigmoid surrogate, $\beta = 25$) on the three Naud parameter sets.
    Top row shows membrane potential over time. Bottom row compares the corresponding spike trains.
    Both channels produce identical voltage traces and identical spike times as expected.}\label{fig:surrogate-channel-verification}
\end{figure}
\clearpage

\begin{figure}[p]
  \begin{center}
    \includegraphics[width=\textwidth]{figures/loss_landscape_focus_vanrossum_a_tau_w_grad2d_superspike_10.pdf}
  \end{center}
  \caption{Van Rossum bound-normalized gradient field on the $a$--$\tau_w$ plane in the \texttt{adaptation} scenario (ground truth $a=2$\,nS, $\tau_w=300$\,ms).
SuperSpike surrogate, $\beta=10$.
The colour scale spans six orders of magnitude in bound-normalized sensitivity, but the quiver is dominated by the horizontal component throughout: $|\partial\mathcal{L}/\partial a \cdot \Delta a| \gg |\partial\mathcal{L}/\partial \tau_w \cdot \Delta \tau_w|$ across the whole slice.}\label{fig:vanrossum-a-tauw}
\end{figure}
\clearpage

\begin{figure}[p]
  \begin{center}
    \includegraphics[width=\textwidth]{figures/subthreshold_fit_quality.pdf}
  \end{center}
  \caption{Fit quality over all subthreshold runs.
  Left: final \gls{MSE} on a clean re-simulation (log scale) --- gradient descent reaches lower loss.
  Right: mean bound-normalized parameter-recovery error --- the two methods are indistinguishable, with gradient descent if anything slightly worse.
  Lower loss does not imply better parameter recovery: the subthreshold loss is degenerate in $\{C_m, E_L, g_L\}$.}\label{fig:subthreshold-fit-quality}
\end{figure}
